% 附录 G 符号、常数与有用公式（Symbols, constants and useful equations）
\chapter{符号、常数与有用公式}

\section{符号表}


\begin{center}
\begin{small}
\begin{tabular}{ll@{\hspace{3em}}ll}
\toprule
$a$ & 单胞尺寸 & $\mathbf{H}$ & 磁场强度\\
$a_0$ & 玻尔半径 & $H_{\mathrm{a}}$ & 外加场\\
$A$ & 面积 & $\mathbf{H}_{\mathrm{d}}$ & 退磁场\\
$A$ & 超精细耦合常数 & $H_{\mathrm{i}}$ & 内场\\
$A$ & 连续介质交换常数 & $\mathcal{H}$ & 哈密顿量\\
$A$ & 质量数 & $\mathrm{i}$ & $\sqrt{-1}$\\
$\mathbf{A}$ & 磁矢势 & $I$ & 电流\\
$b$ & 散射长度 & $I$ & 核自旋量子数\\
$\mathbf{B}$ & 磁通密度（磁场） & $J$ & 总角动量量子数\\
$B_{\mathrm{a}}$ & 外加场 & $\mathsf{J}$ & 交换常数/交换积分\\
$\mathbf{B}_{\mathrm{d}}$ & 退磁场 & $\mathbf{k}$ & 波矢\\
$B_{\mathrm{i}}$ & 内场 & $k_{\mathrm{B}}$ & 玻尔兹曼常数（$1.3807\times10^{-23}$ J\,K$^{-1}$）\\
$\mathbf{B}_{\mathrm{mf}}$ & 分子场 & $k_{\mathrm{F}}$ & 费米波矢\\
$\mathrm{B}_J$ & 布里渊函数 & $K$ & 库仑积分\\
$c$ & 自由空间光速（$2.9979\times10^{8}$ m\,s$^{-1}$） & $K$ & 各向异性常数\\
$\mathbf{D}$ & 电位移 & $K$ & 奈特移位\\
$D$ & 轴各向异性常数 & $K$ & 弹簧常数\\
$\mathrm{d}\mathbf{S}$ & 面元 & $K_{\mathrm{s}}$ & 表面各向异性常数\\
$\mathrm{d}S$ & 熵变 & $K_{\mathrm{v}}$ & 体积各向异性常数\\
$\mathrm{e}$ & $\exp(1)=2.718281828$ & $l$ & 轨道角动量量子数\\
$e$ & 电子电荷大小（$1.6022\times10^{-19}$ C） & $l$ & 朗道能级指标\\
$E$ & 能量 & $m_{\mathrm{e}}$ & 电子质量（$9.109\times10^{-31}$ kg）\\
$E_{\mathrm{F}}$ & 费米能 & $m_l$ & 轨道磁量子数\\
$E_{\mathrm{S}}$ & 单态能量 & $m_I$ & 核磁量子数\\
$E_{\mathrm{T}}$ & 三重态能量 & $m_J$ & 磁量子数\\
$\boldsymbol{\mathcal{E}}$ & 电场 & $m_{\mathrm{n}}$ & 中子质量\\
$F$ & 亥姆霍兹自由能（亥姆霍兹函数） & $m_s$ & 自旋磁量子数\\
$g(E)$ & 能量空间的态密度 & $\mathbf{M}$ & 磁化强度（单位体积磁矩）\\
$g(k)$ & 波矢空间的态密度 & $M_{\pm}$ & 反铁磁体子晶格上的磁化强度\\
$g$ & g 因子 & $M_{\mathrm{s}}$ & 饱和磁化强度\\
$g_{\mathrm{e}}$ & 电子 g 因子 & $n$ & 主量子数\\
$g_I$ & 核 g 因子 & $n$ & 单位体积原子/磁矩数\\
$g_L$ & 轨道角动量 g 因子 & $\mathsf{N}$ & 退磁张量\\
$g_S$ & 自旋角动量 g 因子 & $N$ & 退磁因子\\
$\mathbf{g}$ & 有效 g 张量 & $N$ & 原子核中的中子数\\
$\mathbf{G}$ & 力矩 & $N_{\mathrm{A}}$ & 阿伏伽德罗常数\\
$\mathbf{G}$ & 倒格矢 & $\mathbf{p}$ & 偶极矩\\
$h$ & 普朗克常数（$6.626\times10^{-34}$ J\,s） & $\mathbf{p}$ & 正则动量\\
$\hbar$ & 普朗克常数/$2\pi$（$1.0546\times10^{-34}$ J\,s） & $\hat{\mathbf{p}}$ & 动量算符\\
\bottomrule
\end{tabular}
\end{small}
\end{center}


\begin{center}
\begin{small}
\begin{tabular}{ll@{\hspace{3em}}ll}
\toprule
$\mathbf{q}$ & 磁振子波矢 & $\mu$ & 磁矩\\
$\mathbf{q}$ & 螺旋波矢 & $\mu_{\mathrm{r}}$ & 相对磁导率\\
$\mathbf{Q}$ & 散射矢量 & $\mu_{\mathrm{B}}$ & 玻尔磁子\\
$\mathbf{r}$ & 位置矢量 & $\mu_{\mathrm{N}}$ & 核磁子\\
$\hat{\mathbf{r}}$ & 位置算符 & $\mu_{\mathrm{eff}}$ & 有效磁矩（每化学式单元的 $\mu_{\mathrm{B}}$ 数）\\
$R$ & 电阻 & $\mu_0$ & 自由空间磁导率（$4\pi\times10^{-7}$ H\,m$^{-1}$）\\
$R_{\mathrm{e}}$ & 反常霍尔系数 & $\nu$ & 频率\\
$R_{\mathrm{o}}$ & 正常霍尔系数 & $\xi$ & 关联长度\\
$S$ & 熵 & $\rho$ & 电阻率\\
$s, S$ & 自旋量子数 & $\rho_{\mathrm{H}}$ & 霍尔电阻率\\
$\hat{\mathbf{S}}$ & 自旋角动量算符 & $\boldsymbol{\sigma}$ & 电导率\\
$\hat{S}_{+}$ & 升算符 & $\hat{\boldsymbol{\sigma}}$ & 泡利自旋矩阵\\
$\hat{S}_{-}$ & 降算符 & $\tau$ & 周期\\
$\tilde{S}$ & 有效自旋 & $\tau$ & 超顺磁弛豫时间\\
$t$ & 时间 & $\phi$ & 角度\\
$t$ & 跳跃积分 & $\Phi$ & 磁通量\\
$t_{\mathrm{p}}$ & 脉冲长度 & $\chi$ & 磁化率\\
$T$ & 温度 & $\chi_g$ & 质量磁化率\\
$T_{\mathrm{c}}$ & 临界温度 & $\chi_{\mathrm{L}}$ & 朗道抗磁磁化率\\
$T_{\mathrm{C}}$ & 居里温度 & $\chi_{\mathrm{m}}$ & 摩尔磁化率\\
$T_{\mathrm{N}}$ & 奈尔温度 & $\chi_{\mathrm{P}}$ & 泡利自旋磁化率\\
$T_{\mathrm{P}}$ & Peierls 转变温度 & $\chi_{\mathbf{q}}$ & 波矢依赖磁化率\\
$T_{\mathrm{SP}}$ & 自旋 Peierls 转变温度 & $\chi$ & 自旋波函数\\
$T_1$ & 自旋--晶格弛豫时间 & $\psi$ & 波函数\\
$T_2$ & 自旋--自旋弛豫时间 & $\Psi_{\mathrm{S}}$ & 单态波函数\\
$T_{\mathrm{F}}$ & 费米温度 & $\Psi_{\mathrm{T}}$ & 三重态波函数\\
$U$ & 库仑能 & $\omega$ & 角频率\\
$\mathbf{v}$ & 速度 & $\omega_{\mathrm{c}}$ & 回旋频率\\
$v$ & 速率 & $\omega_{\mathrm{L}}$ & 拉莫尔进动频率\\
$V$ & 势或势能 & $\epsilon_0$ & 自由空间电容率（$8.854\times10^{-12}$ F\,m$^{-1}$）\\
$W$ & 跃迁速率 & $\epsilon_{\mathrm{r}}$ & 相对介电常数\\
$Y_{lm_l}(\theta,\phi)$ & 球谐函数 & $\theta$ & 外斯温度\\
$Z$ & 配分函数 & $\lambda$ & 波长\\
$Z$ & 原子序数 & $\lambda$ & 自旋--轨道常数\\
$\alpha$ & 精细结构常数 & $\lambda_{\text{el-ph}}$ & 电子--声子耦合常数\\
$\beta$ & $=1/k_{\mathrm{B}}T$ & $\lambda_{\text{spin-ph}}$ & 自旋--声子耦合常数\\
$\gamma$ & 旋磁比 & & \\
$\gamma_{\mathrm{e}}$ & 电子旋磁比 & & \\
$\gamma_{\mu}$ & $\mu$ 子旋磁比 & & \\
$\gamma_{\mathrm{N}}$ & 核旋磁比 & & \\
$\Delta$ & 交换劈裂 & & \\
\bottomrule
\end{tabular}
\end{small}
\end{center}

\section{基本常数}

\begin{center}
\smallskip
{基本常数。\\
{\footnotesize\itshape Fundamental constants.}}
\begin{small}
\begin{tabular}{lcl}
\toprule
玻尔半径 & & $a_0 = 5.292\times10^{-11}$ m\\
自由空间光速 & & $c = 2.9979\times10^{8}$ m\,s$^{-1}$\\
电子电荷 & & $e = 1.6022\times10^{-19}$ C\\
普朗克常数 & & $h = 6.626\times10^{-34}$ J\,s\\
 & $h/2\pi=$ & $\hbar = 1.0546\times10^{-34}$ J\,s\\
玻尔兹曼常数 & & $k_{\mathrm{B}} = 1.3807\times10^{-23}$ J\,K$^{-1}$\\
电子静止质量 & & $m_{\mathrm{e}} = 9.109\times10^{-31}$ kg\\
质子静止质量 & & $m_{\mathrm{p}} = 1.6726\times10^{-27}$ kg\\
阿伏伽德罗常数 & & $N_{\mathrm{A}} = 6.022\times10^{23}$ mol$^{-1}$\\
标准摩尔体积 & & $22.414\times10^{-3}$ m$^{3}$\,mol$^{-1}$\\
摩尔气体常数 & & $R = 8.315$ J\,mol$^{-1}$\,K$^{-1}$\\
精细结构常数 & $\dfrac{e^{2}}{4\pi\varepsilon_0\hbar c}=$ & $\alpha = (137.04)^{-1}$\\[8pt]
自由空间电容率 & & $\epsilon_0 = 8.854\times10^{-12}$ F\,m$^{-1}$\\
自由空间磁导率 & & $\mu_0 = 4\pi\times10^{-7}$ H\,m$^{-1}$\\
玻尔磁子 & & $\mu_{\mathrm{B}} = 9.274\times10^{-24}$ A\,m$^{2}$ 或 J\,T$^{-1}$\\
核磁子 & & $\mu_{\mathrm{N}} = 5.051\times10^{-27}$ A\,m$^{2}$ 或 J\,T$^{-1}$\\
中子磁矩 & & $\mu_{\mathrm{n}} = -1.9130\mu_{\mathrm{N}}$\\
质子磁矩 & & $\mu_{\mathrm{p}} = 2.7928\mu_{\mathrm{N}}$\\
\bottomrule
\end{tabular}
\end{small}
\end{center}

\section{有用公式}

（1）恒等式：
\begin{gather*}
\mathrm{e}^{\mathrm{i}\theta} = \cos\theta + \mathrm{i}\sin\theta\\
\sin\theta = \frac{\mathrm{e}^{\mathrm{i}\theta}-\mathrm{e}^{-\mathrm{i}\theta}}{2\mathrm{i}},\qquad
\cos\theta = \frac{\mathrm{e}^{\mathrm{i}\theta}+\mathrm{e}^{-\mathrm{i}\theta}}{2}\\
\sin(\theta+\phi) = \sin\theta\cos\phi+\cos\theta\sin\phi,\qquad
\cos(\theta+\phi) = \cos\theta\cos\phi-\sin\theta\sin\phi\\
\cos^{2}\theta+\sin^{2}\theta = 1\\
\cos2\theta = \cos^{2}\theta-\sin^{2}\theta,\qquad
\sin2\theta = 2\cos\theta\sin\theta\\
\sinh x = \frac{\mathrm{e}^{x}-\mathrm{e}^{-x}}{2},\qquad
\cosh x = \frac{\mathrm{e}^{x}+\mathrm{e}^{-x}}{2}
\end{gather*}

（2）级数展开（$|x|<1$ 时有效）：
\begin{gather*}
(1+x)^{n} = 1+nx+\frac{n(n-1)}{2!}x^{2}+\frac{n(n-1)(n-2)}{3!}x^{3}+\cdots\\
\mathrm{e}^{x} = 1+x+\frac{x^{2}}{2!}+\frac{x^{3}}{3!}+\frac{x^{4}}{4!}+\cdots\\
\sin x = x-\frac{x^{3}}{3!}+\frac{x^{5}}{5!}-\cdots,\qquad
\cos x = 1-\frac{x^{2}}{2!}+\frac{x^{4}}{4!}-\cdots\\
\log(1+x) = x-\frac{x^{2}}{2}+\frac{x^{3}}{3}-\cdots
\end{gather*}

（3）积分：
\begin{gather*}
\int_{0}^{\infty}x^{n}\mathrm{e}^{-x}\,\mathrm{d}x = n!\\
\int_{-\infty}^{\infty}\mathrm{e}^{-\alpha x^{2}}\,\mathrm{d}x = \sqrt{\frac{\pi}{\alpha}}\\
\int_{-\infty}^{\infty}x^{2n}\mathrm{e}^{-\alpha^{2}x^{2}}\,\mathrm{d}x = \frac{(2n)!}{n!2^{2n}\alpha^{n}}\sqrt{\frac{\pi}{\alpha}}
\end{gather*}
不定积分（$a>0$）：
\begin{equation*}
\begin{aligned}
\int\frac{\mathrm{d}x}{x^{2}+a^{2}} &= \frac{1}{a}\tan^{-1}\frac{x}{a}\\
\int\frac{\mathrm{d}x}{x^{2}-a^{2}} &= \frac{1}{2a}\ln\left|\frac{x-a}{x+a}\right|\\
\int\frac{\mathrm{d}x}{\sqrt{x^{2}+a^{2}}} &= \sinh^{-1}\frac{x}{a}\\
\int\frac{\mathrm{d}x}{\sqrt{x^{2}-a^{2}}} &= \begin{cases} \cosh^{-1}\frac{x}{a} & x>a\\ -\cosh^{-1}\frac{x}{a} & x<-a \end{cases}\\
\int\frac{\mathrm{d}x}{\sqrt{a^{2}-x^{2}}} &= \sin^{-1}\frac{x}{a}
\end{aligned}\tag{G.1}
\end{equation*}

（4）矢量算符：
\begin{itemize}[itemsep=0.2em]
\item grad（梯度）作用于标量场产生矢量场：
$\operatorname{grad}\psi=\nabla\psi=\left(\frac{\partial\psi}{\partial x},\frac{\partial\psi}{\partial y},\frac{\partial\psi}{\partial z}\right)$
\item div（散度）作用于矢量场产生标量场：
$\operatorname{div}\mathbf{A}=\nabla\cdot\mathbf{A}=\frac{\partial A_x}{\partial x}+\frac{\partial A_y}{\partial y}+\frac{\partial A_z}{\partial z}$
\item curl（旋度）作用于矢量场产生另一个矢量场：
$\operatorname{curl}\mathbf{A}=\nabla\times\mathbf{A}=\begin{vmatrix}\mathbf{i} & \mathbf{j} & \mathbf{k}\\ \partial/\partial x & \partial/\partial y & \partial/\partial z\\ A_x & A_y & A_z\end{vmatrix}$
\end{itemize}
其中 $\psi(\mathbf{r})$ 与 $\mathbf{A}(\mathbf{r})$ 分别是任意给定的标量场与矢量场。

（5）矢量恒等式：
\begin{align*}
&\nabla\cdot(\nabla\psi) = \nabla^{2}\psi\\
&\nabla\times(\nabla\psi) = 0\\
&\nabla\cdot(\nabla\times\mathbf{A}) = 0\\
&\nabla\cdot(\psi\mathbf{A}) = \mathbf{A}\cdot\nabla\psi+\psi\nabla\cdot\mathbf{A}\\
&\nabla\times(\psi\mathbf{A}) = \psi\nabla\times\mathbf{A}-\mathbf{A}\times\nabla\psi\\
&\nabla\times(\nabla\times\mathbf{A}) = \nabla(\nabla\cdot\mathbf{A})-\nabla^{2}\mathbf{A}\\
&\nabla\cdot(\mathbf{A}\times\mathbf{B}) = \mathbf{B}\cdot\nabla\times\mathbf{A}-\mathbf{A}\cdot\nabla\times\mathbf{B}\\
&\nabla(\mathbf{A}\cdot\mathbf{B}) = (\mathbf{A}\cdot\nabla)\mathbf{B}+(\mathbf{B}\cdot\nabla)\mathbf{A}+\mathbf{A}\times(\nabla\times\mathbf{B})+\mathbf{B}\times(\nabla\times\mathbf{A})\\
&\nabla\times(\mathbf{A}\times\mathbf{B}) = (\mathbf{B}\cdot\nabla)\mathbf{A}-(\mathbf{A}\cdot\nabla)\mathbf{B}+\mathbf{A}(\nabla\cdot\mathbf{B})-\mathbf{B}(\nabla\cdot\mathbf{A})
\end{align*}
这些恒等式用交错张量与求和约定即可轻松证明。交错张量 $\epsilon_{ijk}$ 定义为
\begin{equation*}
\epsilon_{ijk} = \begin{cases} 1 & \text{若 } ijk \text{ 是 } 123 \text{ 的偶置换}\\ -1 & \text{若 } ijk \text{ 是 } 123 \text{ 的奇置换}\\ 0 & \text{若 } i, j, k \text{ 中任何两个相等} \end{cases}
\end{equation*}
于是矢量积可以写作
\begin{equation*}
(\mathbf{A}\times\mathbf{B})_i = \epsilon_{ijk}A_jB_k.
\end{equation*}
这里使用求和约定：重复两次的指标默认求和。标量积于是为
\begin{equation*}
\mathbf{A}\cdot\mathbf{B} = A_iB_i.
\end{equation*}
可以利用恒等式
\begin{equation*}
\epsilon_{ijk}\epsilon_{ilm} = \delta_{jl}\delta_{km}-\delta_{jm}\delta_{kl}
\end{equation*}
其中 $\delta_{ij}$ 是克罗内克 delta：
\begin{equation*}
\delta_{ij} = \begin{cases} 1 & i=j\\ 0 & i\neq j \end{cases}
\end{equation*}
矢量三重积为
\begin{equation*}
\mathbf{A}\times(\mathbf{B}\times\mathbf{C}) = (\mathbf{A}\cdot\mathbf{C})\mathbf{B}-(\mathbf{A}\cdot\mathbf{B})\mathbf{C}.
\end{equation*}

（6）柱坐标：
\begin{gather*}
\nabla^{2}\psi = \frac{1}{r}\frac{\partial}{\partial r}\left(r\frac{\partial\psi}{\partial r}\right) + \frac{1}{r^{2}}\frac{\partial^{2}\psi}{\partial\phi^{2}} + \frac{\partial^{2}\psi}{\partial z^{2}}\\
\nabla\psi = \left(\frac{\partial\psi}{\partial r}, \frac{1}{r}\frac{\partial\psi}{\partial\phi}, \frac{\partial\psi}{\partial z}\right)
\end{gather*}

（7）球极坐标：
\begin{align*}
\nabla^{2}\psi &= \frac{1}{r^{2}}\frac{\partial}{\partial r}\left(r^{2}\frac{\partial\psi}{\partial r}\right) + \frac{1}{r^{2}\sin\theta}\frac{\partial}{\partial\theta}\left(\sin\theta\frac{\partial\psi}{\partial\theta}\right) + \frac{1}{r^{2}\sin^{2}\theta}\frac{\partial^{2}\psi}{\partial\phi^{2}}\\
\nabla\psi &= \left(\frac{\partial\psi}{\partial r}, \frac{1}{r}\frac{\partial\psi}{\partial\theta}, \frac{1}{r\sin\theta}\frac{\partial\psi}{\partial\phi}\right)
\end{align*}

\vspace{2em}
\begin{center}
\small\kaishu ——全书完——
\end{center}
